\usepackage{enumitem}
\usepackage[utf8]{inputenc} % allow utf-8 input
\usepackage[T1]{fontenc}    % use 8-bit T1 fonts
\usepackage{hyperref}       % hyperlinks
\usepackage{url}            % simple URL typesetting
\usepackage{mathrsfs}
\usepackage{booktabs}       % professional-quality tables
\usepackage{amsfonts}       % blackboard math symbols
\usepackage{nicefrac}       % compact symbols for 1/2, etc.
\usepackage{microtype}      % microtypography
\usepackage{xcolor}         % colors
\usepackage{xcolor}


\usepackage{microtype}
\usepackage{graphicx}
\usepackage{subfigure}
\usepackage{booktabs} % for professional tables
\usepackage{bbm}
\newcommand\R{\mathbb{R}}
% \newcommand{\ishani}[1]{\textcolor{blue}{#1}}
% \newcommand{\ishanicomment}[1]{\textcolor{red}{#1}}

% hyperref makes hyperlinks in the resulting PDF.
% If your build breaks (sometimes temporarily if a hyperlink spans a page)
% please comment out the following usepackage line and replace
% \usepackage{icml2025} with \usepackage[nohyperref]{icml2025} above.
\usepackage{hyperref}

% Attempt to make hyperref and algorithmic work together better:
\usepackage{algorithm}
\usepackage[noend]{algorithmic}
\newcommand{\Comment}[1]{\textcolor{gray}{\footnotesize\textit{// #1}}}


% Use the following line for the initial blind version submitted for review:

% If accepted, instead use the following line for the camera-ready submission:
% \usepackage[accepted]{icml2025}

% For theorems and such
\usepackage{amsmath}
\usepackage{amssymb}
\usepackage{mathtools}
\usepackage{amsthm}

% if you use cleveref..
\usepackage[capitalize,noabbrev]{cleveref}

\usepackage{tikz}
\usepackage{tikz-cd}
\usetikzlibrary{automata,positioning}

%%%%%%%%%%%%%%%%%%%%%%%%%%%%%%%%
% THEOREMS
%%%%%%%%%%%%%%%%%%%%%%%%%%%%%%%%
\theoremstyle{plain}
\newtheorem{theorem}{Theorem}[section]
\newtheorem{proposition}[theorem]{Proposition}
\newtheorem{example}[theorem]{Example}

\newtheorem{lemma}[theorem]{Lemma}
\newtheorem{corollary}[theorem]{Corollary}
\theoremstyle{definition}
\newtheorem{definition}[theorem]{Definition}
\newtheorem{assumption}[theorem]{Assumption}
\newtheorem{remark}[theorem]{Remark}
\newtheorem{discussion}[theorem]{Discussion}
\newtheorem{property}[theorem]{Property}
\def\ceil#1{\left\lceil #1 \right\rceil}

\def\est{\mathrm{est}}
\def\E{\mathbb{E}}
\def\N{\mathbb{N}}
\DeclareMathOperator*{\argmax}{arg\,max}

\newcommand{\Real}{\mathbb{R}}
\newcommand{\Nat}{\mathbb{N}}
\newcommand{\F}{\mathbb{F}}
\newcommand{\Z}{\mathbb{Z}}
\renewcommand{\S}{\mathbb{S}}%

% probability
\renewcommand{\P}{\mathbb{P}}

\DeclareMathOperator*{\argmin}{argmin}
\newcommand{\poly}{\operatorname{poly}}
\newcommand{\polylog}{\operatorname{polylog}}
\newcommand{\mdiag}{\mathbf{Diag}} % diagonal matrix with vector on diag
\newcommand{\sign}{\operatorname{sign}}%
\newcommand{\nnz}{\operatorname{nnz}}%
\newcommand{\tail}{\operatorname{tail}}%
\newcommand{\dist}{\operatorname{dist}}%

% redfine some classical notation
\renewcommand{\tilde}{\widetilde}
\renewcommand{\hat}{\widehat}
\renewcommand{\bar}{\overline}

\usepackage{pgffor}
\foreach \x in {A,...,Z}{
		\expandafter\xdef\csname m\x\endcsname{\noexpand\mathbf{\x}}
	}
\foreach \x in {A,...,Z}{
		\expandafter\xdef\csname om\x\endcsname{\noexpand\overline{\noexpand\mathbf{\x}}}
	}
\foreach \x in {A,...,Z}{
		\expandafter\xdef\csname c\x\endcsname{\noexpand\mathcal{\x}}
	}


\makeatletter

\DeclareRobustCommand\widecheck[1]{{\mathpalette\@widecheck{#1}}}
\def\@widecheck#1#2{%
	\setbox\z@\hbox{\m@th$#1#2$}%
	\setbox\tw@\hbox{\m@th$#1%
			\widehat{%
				\vrule\@width\z@\@height\ht\z@
				\vrule\@height\z@\@width\wd\z@}$}%
	\dp\tw@-\ht\z@
	\@tempdima\ht\z@ \advance\@tempdima2\ht\tw@ \divide\@tempdima\thr@@
	\setbox\tw@\hbox{%
		\raise\@tempdima\hbox{\scalebox{1}[-1]{\lower\@tempdima\box
				\tw@}}}%
	{\ooalign{\box\tw@ \cr \box\z@}}}


\newcommand{\manifold}{\mathcal{M}}
\newcommand{\dataset}{\mathcal{D}}
\newcommand{\domain}{\mathcal{X}}
\newcommand{\ffam}{\mathcal{F}}
\newcommand{\relu}{\mathrm{ReLU}}
\newcommand{\vspan}{\mathrm{span}}

\newcommand{\todo}[1]{\textcolor{red}{[TODO: #1]}}
\newcommand{\nastia}[1]{\textcolor{pink}{[Nastia says: #1]}}
\newcommand{\emile}[1]{\textcolor{blue}{[Emile says: #1]}}
\newcommand{\pan}[1]{\textcolor{magenta}{[Pan says: #1]}}
\newcommand{\wenjing}[1]{\textcolor{pink}{[Wenjing says: #1]}}

\newcommand{\TV}{\mathrm{TV}} % total variation norm
\newcommand{\deff}{d_{\mathrm{eff}}} % effective dimension
\newcommand{\CD}{\mathrm{CD}} % curvature-dimension condition

\newcommand{\kernel}{\kappa}

\newcommand{\HK}{\mathcal{H}_\kernel} % Hilbert space of functions induced by a kernel $K$


% \newcommand{\HK}{\mathcal{H}_K}
\newcommand{\X}{\mathcal{X}}
\newcommand{\B}{\mathcal{B}}
% \newcommand{\E}{\mathbb{E}}
% \newcommand{\R}{\mathbb{R}}
% \newcommand{\N}{\mathbb{N}}
\newcommand{\iid}{\stackrel{\mathrm{iid}}{\sim}}
\newcommand{\inprod}[2]{\langle #1, #2 \rangle}
\newcommand{\norm}[1]{\| #1 \|}

\DeclareMathOperator{\Cov}{Cov}
\DeclareMathOperator{\Var}{Var}
\DeclareMathOperator{\Range}{Range}
\DeclareMathOperator{\Tr}{Tr}
\newcommand{\Splus}{\mathcal{S}_+^d}
\newcommand{\Splusplus}{\mathcal{S}_{++}^d}



